\documentclass[xcolor=svgnames,11pt]{beamer}

\usepackage{graphicx}

\usepackage[italian, english]{babel}
\usepackage[utf8x]{inputenc}
\usepackage[T1]{fontenc}
\usepackage[scale=0.9,sfdefault,light]{roboto}

\usepackage{hyperref}
\usepackage{tabularx}
\usepackage{url}
\usepackage{pbox}

\usetheme{Goettingen}
\definecolor{leorange}{HTML}{FBA100}
\definecolor{leblue}{HTML}{003A70}
\definecolor{myblack}{RGB}{14,14,14}
\definecolor{lelightblue}{HTML}{BFCDDB}

\usecolortheme[named=leblue]{structure}
\beamertemplatenavigationsymbolsempty

\beamertemplateshadingbackground{white!15}{leorange!5}
\setbeamertemplate{sidebar canvas right}[vertical shading]%
[top=lelightblue,bottom=lelightblue]

\mode<presentation>{\usebackgroundtemplate{\includegraphics
    [height=\paperheight]{img/back_le.png}}}

\setbeamertemplate{blocks}[rounded][shadow=false]
\setbeamercolor{block body}{bg=leorange!50}
\setbeamercolor{block title}{bg=leorange!80}

\setbeamercovered{transparent}

\title{\textbf{Let's encrypt! Certificati SSL per tutti!}}
\author{Nicola Corti}
\institute{Gruppo Utenti Linux Pisa \\ \medskip \includegraphics[height=1.5cm]{img/gulp_logo.png}}
\logo{img/gulp_logo.png}
\date{02 Febbraio 2016}


\begin{document}

\begin{frame}
	\titlepage
\end{frame}

\section{Intro}

\begin{frame}{Cosa è Let's encrypt?}

\begin{center}
\includegraphics[width=5cm]{img/logo_le.png}
\end{center}

\begin{block}{}
Let's encrypt è una \textbf{Certification Authority (CA)} che rilascia
certificati SSL/TLS a titolo \textbf{gratuito}.
\end{block}
\pause \medskip
\begin{itemize}
  \item Dal 5 Dicembre 2015 si trova in \textbf{Public Beta}.
  \pause
  \item Ad oggi ha rilasciato più di \textbf{480 k} certificati.
  \pause
  \item Fa dell'\textbf{automazione} il suo punto di forza.
  \pause
  \item \underline{https://letsencrypt.org/}
\end{itemize}
\end{frame}

\begin{frame}{Sponsors}
\begin{center}
Chi sta supportando il progetto?
\medskip
\vspace{0.5cm}
\includegraphics[width=10cm]{img/sponsors.png}
\end{center}
\end{frame}

\begin{frame}{History}
\begin{description}
  \item[2012] Iniziato il progetto dentro \textbf{Mozilla}
  \pause
  \item[11-2014] Let's Encrypt \textbf{annunciato pubblicamente}
  \pause
  \item[01-2015] Protocollo \textbf{ACME} sottoposto all'\textbf{IETF} per essere standardizzato.
  \pause
  \item[04-2015] \textbf{ISRG} (Internet Security Research Group) e la \textbf{Linux Foundation}
  si uniscono al progetto
  \pause
  \item[09-2015] Rilasciato il primo certificato per \textbf{helloworld.letsencrypt.org}
  \pause
  \item[10-2015] Il certificato intermedio di L.E. diventa \textbf{\emph{cross-signed}} grazie a IdenTrust.
  (\emph{Let's Encrypt is trusted!})
  \pause
  \item[12-2015] \textbf{Public beta!}
\end{description}
\end{frame}

\begin{frame}{Why?}
\begin{itemize}
  \item \textbf{Free}
  \item \textbf{Automatico}
  \item \textbf{Open}
  \item \textbf{Sicuro}
  \item \textbf{Trasparente}
\end{itemize}
\end{frame}

\section{Server}
\begin{frame}{Boulder}
\begin{block}{Architettura}
Il sistema di Let's Encrypt si basa su \textbf{3 Componenti}: un \textbf{server},
un \textbf{client} e il \textbf{protocollo} che permette a server e client di comunicare.
\end{block}

\medskip \pause

Il server si chiama \textbf{Boulder} ed è scritto interamente in \textbf{Go}.
Si occupa di gestire tutte le procedure per il \textbf{rilascio} e \textbf{revoca} dei certificati.
\pause
\begin{center}
\visible<3>{\includegraphics[width=0.3cm]{img/logo_gh.pdf}}\textbf{ https://github.com/letsencrypt/boulder}
\end{center}
\end{frame}



\section{Client}
\section{Protocol}
\section{Certs}
\section{Setup}

\begin{frame}{}
\begin{center}
\begin{Huge}
{\color{leorange} \textbf{Domande...?}}
\end{Huge}

\vspace{1.5cm}
\begin{small}
Slides realizzate da:\\
\textbf{Nicola Corti - corti.nico [at] gmail [dot] com}\\

\bigskip

Slides realizzate con \LaTeX\ Beamer.\\
La seguente presentazione \`e rilasciata sotto licenza\\
\begin{footnotesize}	\textbf{Creative Commons - Attributions, Non Commercial, Share-alike}.
\end{footnotesize}
\\
\medskip
\includegraphics[height=0.5cm]{img/cc.png}

\end{small}
\end{center}
\end{frame}

\end{document}
